% Appendix A -- parameter measurement, moved out of section 3 (2026-09-26) to
% keep the body within the MLSys page limit. Nothing here is needed to follow
% the argument; all of it is needed to check it.

\section{Measured parameters and sensitivity}
\label{app:params}

A100-80 PCIe; two models; 41 single-shot configurations plus a
13-configuration $n{=}3$ repeat sweep; distinct random contexts; prefix caching
disabled.
% src: paper/model/p4_llama8b_a100_rep3.json @ 2bfc6ca5e

\paragraph{Weight read is efficient and stable within a model.} $\effw \approx
0.75$ of the measured $1935$\,GB/s peak ($\approx$1450\,GB/s achieved). The
$B{=}1$ steps must have their KV term subtracted before $\effw$ is read off
them: it is 1.6\% of the step at 2K context and 6.0\% at 8K, and leaving it in
manufactures a drift (0.729 falling to 0.694 for Llama-3.1-8B) that is the KV
term growing rather than the weight read degrading. Subtracted, $\effw$ is
0.741/0.736/0.738 for Llama-3.1-8B across 2K--8K and 0.778/0.777 for
Qwen2.5-32B --- flat within each model, and 5.3\% apart between them. We carry
the single value 0.75, which is between them, and \S\ref{sec:model:cancel}
reports what each choice predicts.

\paragraph{Paged KV read is neither.} $\effkv$ degrades by $1.6\times$ with the
KV working set (Table~\ref{tab:effkv}):

\begin{table*}[t]
\centering
\small
\begin{tabular}{lrrrrr}
\toprule
KV working set & 4.3\,GB & 8.6\,GB & 17--34\,GB & 43\,GB & 51.5\,GB \\
\midrule
$\effkv$ & 0.77 & 0.72 & 0.54--0.56 & 0.49 & 0.46 \\
\bottomrule
\end{tabular}
\caption{Paged KV read efficiency is not a constant. Error bars from $n{=}3$
repeats are $\pm0.02$--$0.04$; the degradation is far outside them. The
17--34\,GB column takes 0.55 rather than the 0.747 measured at
$B{=}16,S{=}8192$ --- the anomaly discussed below. That is the
\emph{optimistic} choice for us, since a lower $\effkv$ lengthens the step and
so raises the projected gain, and we note it rather than let it pass.}
\label{tab:effkv}
\end{table*}

This is the real correction to the model. It also bears on how such models are
built elsewhere: \textbf{a roofline that charges KV at one fixed bandwidth will
be optimistic at the capacity-bound operating point}, because that is where
$\effkv$ is worst. We measured this on one GPU and one kernel generation
(A100, FlashAttention-2 era), so we state it as a property of what we measured
rather than of the literature.

\paragraph{A caveat at kernel granularity.} The claim that only total KV bytes
matter holds at three of four same-byte test points but breaks at 17.2\,GB
($B{=}32,S{=}4096$ gives 0.561 against $B{=}16,S{=}8192$ at 0.747). The cause
is a reproducible anomaly at $B{=}32,S{=}4096$ that also appeared in an
independent earlier run --- an attention-kernel scheduling artifact, most
likely a split-KV heuristic in the attention kernel choosing a different
partitioning at the larger batch, not measurement noise ($\pm0.61$\,ms
against a $\sim$3\,ms deviation). The aggregate law holds; the
per-configuration surface has kernel-regime steps.

\paragraph{Compute is not the binding term.} Decode-attention arithmetic
intensity equals the GQA group size --- 8 for 70B, 4 for 8B --- against an
H100 balance point near 116 FLOP/byte, so attention sits 14--29$\times$ below
the compute roof. At batch 8, 70B spends 7\% of its step on arithmetic.
% src: paper/model/compute_check.py @ 2bfc6ca5e
Llama-8B at 4K crosses the arithmetic-intensity balance point at batch 142, while its step remains memory-bound by about $4\times$, which is a further
reason it is the wrong model to evaluate on.



\subsection*{The step-time fits in full}

\paragraph{Validated on hardware, per model.} Equation~\ref{eq:step} makes two
claims: that ITL depends on batch and context only through their product
$B\cdot S$, and that the intercept is the weight term $W/(\BW\effw)$. We test
them on an A100-80 over 41 single-shot configurations --- 27 on Llama-3.1-8B, 14
on Qwen2.5-32B --- plus 13 three-repeat points drawn from the same grid.

Within each model, ITL is linear in $B\cdot S$: $R^2 = 0.975$ for Llama-3.1-8B
with residuals unbiased across $S \in \{2048, 4096, 8192\}$ (means $+0.09$,
$-0.09$, $+0.00$\,ms), and $R^2 = 0.92$ for Qwen2.5-32B. Unbiased in $S$ is the
weaker statement it sounds like: grouped by KV \emph{magnitude} the same
residuals are a clean U, from $+2$\,ms below 1\,GB through $-6.6$ in the middle
to $+6$\,ms (8\%) at the largest points. That curvature is not noise --- it is
the $\effkv$ degradation of \S\ref{sec:model:params} showing through a linear
fit, which is why we use the measured $\effkv$ curve rather than the regression
for anything that matters. $B{=}16,S{=}4096$ and
$B{=}8,S{=}8192$ both measure $17.17$\,ms. \textbf{Pooling the two models is
meaningless and we do not report it} --- the fit across both is $R^2 = 0.19$,
because Eq.~\ref{eq:step}'s intercept is model-dependent, which is the whole
content of \S\ref{sec:model:modeldep}.

That intercept is the sharper test, because Eq.~\ref{eq:step} predicts it rather
than fitting it. Predicting it needs $\effw$, which we extract from the $B{=}1$
points after subtracting their KV term --- small but not negligible, 1.6\% of
the step at 2K context and 6.0\% at 8K, and leaving it in produces a spurious
drift in $\effw$ with context. Subtracted, $\effw$ is stable within each model
(0.741, 0.736, 0.738 for Llama-3.1-8B across 2K--8K; 0.778, 0.777 for
Qwen2.5-32B) and differs 5.3\% between them.

Against those values Eq.~\ref{eq:step} predicts intercepts of $11.2$ and
$43.6$\,ms where the regressions fit $9.4$ and $43.6$. \textbf{The decomposition
is exact for Qwen2.5-32B --- within $0.01$\,ms --- and wrong by $1.9$\,ms, or
20\%, for Llama-3.1-8B.} We report the split rather than the average because the
average would hide it: a single $\effw = 0.75$ gives $11.1$ and $45.2$\,ms, both
overshooting, which looks like uniform miscalibration and is not.

The 8B gap is the model's clearest failure on the data we have. Its most likely
source is that Eq.~\ref{eq:step} charges weights and KV at efficiencies fitted
in separate regimes, and an 8B's step is 86\% KV (\S\ref{sec:background}), so
almost none of it is the term the intercept isolates --- which is also why the
8B is the row the projections should be trusted least on.
% src: paper/model/p4_llama8b_a100.json, p4_qwen32b_a100.json @ 2bfc6ca5e
We use differential timing --- $\mathrm{ITL} = [T(N_2) - T(N_1)]/(N_2 - N_1)$
with identical prompts in both runs, so prefill and warmup cancel; context
grows across the $N_2-N_1$ tokens, which the differencing does not remove.


\subsection*{Sensitivity, and what is held out}

\paragraph{Weight precision is the clean lever on $\HBM/\Ckv$, and we do not
pull it.} Comparing models of different sizes moves the weight footprint, but
also layer count, KV-head count and architecture, so a trend across 8B, 32B and
70B confounds the quantity the law is about with three others. The clean
instrument is a \emph{single} model at two weight precisions, which moves
$\HBM/\Ckv$ and nothing else --- the two Qwen-32B rows above share an
architecture, a tokenizer and a KV geometry and differ only in that fp8 weights
leave 60.0\,GB for KV where fp16 leave 26.5\,GB.

We report those rows as projections rather than measurements. Validating the
$\HBM/\Ckv$ axis directly would have cost hours we judged better spent
comparing against the closest prior system, so the law is validated across
\emph{bandwidth} (\S\ref{sec:eval:cost}) while its amplifier rests on the
parameter fit and on agreement with published results
(\S\ref{sec:model:directkv}). \S\ref{sec:eval:model} states this again where
a reviewer will look for it.

The same observation has a practical edge, which \S\ref{sec:discussion} takes
up: quantizing \emph{weights} is not neutral to a memory system, it actively
removes its headroom, while quantizing \emph{KV} scales both cohorts together
and composes.

\paragraph{What is and is not held out.} $\effw$ is calibrated on the $B{=}1$
points, with their KV term subtracted, and $\effkv$ on sub-capacity working
sets, so the two come from disjoint parts of the same sweeps. The full-pool step
times above are \emph{not} an independent check on them --- they are computed
from those same fits, so comparing Eq.~\ref{eq:step} against them would be
circular and we do not. The one out-of-sample test we have is the intercept
recovery of \S\ref{sec:model:cancel}, and it passes on one model of two:
exactly on Qwen2.5-32B, and 20\% high on Llama-3.1-8B. Measuring a real
full-pool step on a device large enough to hold one would be the stronger test,
and we have not run it.

Several of these parameters are measured with error, so it is worth asking
which ones the projection actually depends on. For Qwen-32B fp16 on GH200-96 at
8K --- where Table~\ref{tab:gain} reads $\times1.45$ --- $\effw$ over 0.70--0.80
barely moves it, spanning $\times1.41$--$\times1.47$. \textbf{$\effkv$ is a
different matter}: $\pm20\%$ spans $\times1.31$--$\times1.47$, and the downside
is the larger half. Charging the
rider's HBM passes is what made it so, because $\effkv$ now prices the rider
traffic as well as the resident cohort's, and it is the one parameter a new
kernel generation moves. The reason is structural --- the gain depends on a
\emph{ratio}, so a bandwidth error largely cancels between the step-time
numerator and the capacity denominator.

The other parameter that matters is the \textbf{weight footprint}. Moving it
from 60 to 68\,GB spans $\times1.34$--$\times1.53$, because weights sit in
$\Ckv = \HBM - W$ and are cancelled by nothing. This is why we use measured
checkpoint sizes rather than nominal parameter counts --- for a quantized model
the two differ by gigabytes and the error propagates undamped into the headline
--- and it is why halving weight precision is the sharpest lever we have on
$\HBM/\Ckv$.
% src: paper/model/killer_v2.py sensitivity sweep, 2026-09-06

Note the direction of the $\effkv$ term, because it is our main downside risk:
a \emph{more} efficient paged-KV read shortens the step, shrinks the transfer
budget, and \emph{lowers} the gain. The curve above is measured on FA-2-era
kernels, and a Hopper-generation kernel~\cite{flashattention3} reads paged KV
differently. Re-measuring it there is a prerequisite for any Hopper
projection, and a downside risk rather than a formality: a kernel generation
that reads paged KV \emph{more} efficiently shortens the step, shrinks the
transfer budget, and lowers the predicted gain.
